% The Slide Definitions
%document
\documentclass[10pt]{beamer}
%theme
\usetheme{metropolis}
% packages
\usepackage{color}
\usepackage{listings}
\usepackage[ngerman]{babel}
\usepackage[utf8]{inputenc}
\usepackage{multicol}


% color definitions
\definecolor{mygreen}{rgb}{0,0.6,0}
\definecolor{mygray}{rgb}{0.5,0.5,0.5}
\definecolor{mymauve}{rgb}{0.58,0,0.82}

\lstset{
    backgroundcolor=\color{white},
    % choose the background color;
    % you must add \usepackage{color} or \usepackage{xcolor}
    basicstyle=\footnotesize\ttfamily,
    % the size of the fonts that are used for the code
    breakatwhitespace=false,
    % sets if automatic breaks should only happen at whitespace
    breaklines=true,                 % sets automatic line breaking
    captionpos=b,                    % sets the caption-position to bottom
    commentstyle=\color{mygreen},    % comment style
    % deletekeywords={...},
    % if you want to delete keywords from the given language
    extendedchars=true,
    % lets you use non-ASCII characters;
    % for 8-bits encodings only, does not work with UTF-8
    frame=single,                    % adds a frame around the code
    keepspaces=true,
    % keeps spaces in text,
    % useful for keeping indentation of code
    % (possibly needs columns=flexible)
    keywordstyle=\color{blue},       % keyword style
    % morekeywords={*,...},
    % if you want to add more keywords to the set
    numbers=left,
    % where to put the line-numbers; possible values are (none, left, right)
    numbersep=5pt,
    % how far the line-numbers are from the code
    numberstyle=\tiny\color{mygray},
    % the style that is used for the line-numbers
    rulecolor=\color{black},
    % if not set, the frame-color may be changed on line-breaks
    % within not-black text (e.g. comments (green here))
    stepnumber=1,
    % the step between two line-numbers.
    % If it's 1, each line will be numbered
    stringstyle=\color{mymauve},     % string literal style
    tabsize=4,                       % sets default tabsize to 4 spaces
    % show the filename of files included with \lstinputlisting;
    % also try caption instead of title
    language = [Sharp]C,
	showspaces = false,
	showtabs = false,
	showstringspaces = false,
	escapechar = ,
}

\def\ContinueLineNumber{\lstset{firstnumber=last}}
\def\StartLineAt#1{\lstset{firstnumber=#1}}
\let\numberLineAt\StartLineAt



\newcommand{\codeline}[1]{
	\alert{\texttt{#1}}
}


% Author and Course information
% This Document contains the information about this course.

% Authors of the slides
\author{\href{https://github.com/evilham}{@evilham}\\
\footnotesize{Basierend auf Folien von \href{https://github.com/satkowski}{@satkowski} (November 2016)}}

% Name of the Course
\institute{C\texttt{\#} Kurs}

% Fancy Logo
\titlegraphic{\hfill\includegraphics[height=1.25cm]{../templates/fsr_logo_cropped}}


% Presentation title
% TODO Change the topic of the lesson
\title{Awesome Stuff}
\date{\today}


\begin{document}

\maketitle

\begin{frame}{Gliederung}
	\setbeamertemplate{section in toc}[sections numbered]
	\tableofcontents
\end{frame}

% ----------------------- null ------------------------------
\section{null}
\begin{frame}{null}
	\begin{itemize}
		\item Schlüsselwort für eine Referenz die auf kein Objekt zeigt
		\item Alle Referenztyp-Variablen können \alert{null} sein
		\item Mit/auf \alert{null} kann man nicht arbeiten, sonst gibt es eine ArgumentNullException
		\begin{itemize}
			\item Außer ein Vergleich ob eine Variable \alert{null} ist
		\end{itemize}
	\end{itemize}
	\lstinputlisting{resources/05_null_exceptions_schluesselwoerter/null.cs}	
\end{frame}

\subsection{Nullable}
\begin{frame}{Nullable}
	\begin{itemize}
		\item Werden genutzt, damit Werttyp-Variablen ebenfalls \alert{null} haben können
		\item Besitzt 2 Eigenschaften
		\begin{description}
			\item[HasValue] Gibt an, ob das Objekt einen Wert hat oder nicht
			\item[Value] Gibt den Wert des Objektes zurück\\ Wirft eine InvalidOperationException falls HasValue false ist 
		\end{description}
		\item Syntax:
		\begin{itemize}
			\item \texttt{\alert{<Werttyp>}? \alert{<Variablenname>};} oder
			\item \texttt{Nullable<\alert{<Datentyp>}> \alert{<Variablenname>};}
			\begin{itemize}
				\item Das ist ein Generic
			\end{itemize}
		\end{itemize}
	\end{itemize}
\end{frame}

\begin{frame}{Nullable}
	\lstinputlisting{resources/05_null_exceptions_schluesselwoerter/nullable.cs}	
\end{frame}

% ----------------------- Exceptions ------------------------------
\section{Exceptions}
\begin{frame}{Exceptions}
	\begin{itemize}
		\item Zu dt. Ausnahmen
		\item Zeigen an wenn im Programm eine Situation aufgetreten ist, die es nicht lösen kann
		\item Beinhalten Information über Art und Auftreten des Fehlers
		\item Es gibt verschiedene Arten, wie z.B.:
		\begin{itemize} 
			\item ArgumentException
			\item ArgumentNullException	
			\item FileNotFoundException
		\end{itemize}
		\item Sie bringen das Programm zum Abstürzen, wenn diese nicht vom Code behandelt werden
	\end{itemize}
\end{frame}

\subsection{Handling}
\begin{frame}{Handling}
	\textbf{try catch finally}\\
	\begin{itemize}
		\item Wird genutzt um Exceptions aufzufangen und diese zu behandeln
		\item Der im \alert{try}-Block stehende Code wird ausgeführt und bei auftreten einer Exception werden die \alert{catch}-Blöcke ausgeführt
		\item Es kann mehrere \alert{catch}-Blöcke für verschiedene Exception geben
		\begin{itemize}
			\item Die Exception kann im \alert{catch}-Block als Variable genutzt werde
			\item Man geht immer vom spezifischen zum allgemeinem Fehler, damit alle Exceptions gefangen werden können
		\end{itemize}
		\item Der \alert{finally}-Block wird immer ausgeführt, auch wenn keine Exception aufgetreten ist
		\begin{itemize}
			\item Wird auch Bereinigungsblock genannt
		\end{itemize}
	\end{itemize}
\end{frame}

\begin{frame}{Handling}
	\lstinputlisting{resources/05_null_exceptions_schluesselwoerter/exception_handling.cs}	
\end{frame}

\subsection{Eigene Exceptions}
\begin{frame}{Eigene Exceptions}
	\textbf{throw}\\
	\begin{itemize}
		\item Schlüsselwort um eine Exception zu werfen/auszulösen
		\item Syntax:
		\begin{itemize}
			\item \texttt{throw new \alert{<Exceptionkonstruktor>}(\alert{[Parameter]});}
		\end{itemize}
	\end{itemize}
	\textbf{Eigene Exception}\\
	\begin{itemize}
		\item Man kann eine eigene Excpetion-Klasse schreiben, in dem man von der Klasse \alert{System.Exception} erbt
		\item Dabei ist darauf zu achten, dass man gegeben falls die geerbten Konstruktoren implementiert
	\end{itemize}
\end{frame}

\begin{frame}{Eigene Exceptions}
	\lstinputlisting{resources/05_null_exceptions_schluesselwoerter/exception_throw.cs}
\end{frame}

% ----------------------- Schlüsselwörter ------------------------------
\section{Schlüsselwörter}
\begin{frame}{Schlüsselwörter}
	\textbf{static}\\
	\begin{itemize}
		\item Es können Klassen, Methoden, Eigenschaften, Felder und mehr statisch sein
		\item Damit gehören Felder und co. nicht mehr zu einen Objekt einer Klasse sondern zu dem Typ selbst 
		\item Werden über die Klasse/Datentyp aufgerufen
		\item Methoden können nicht auf nicht-statische Felder/Attribute der eigenen Klasse zugreifen
	\end{itemize}
	\textbf{readonly}\\
	\begin{itemize}
		\item Kann nur auf Felder benutzt werden
		\item Stellt sicher, dass das Feld im Laufe der Benutzung nicht verändert wird
		\item Der Wert des Feldes kann nur durch direkte Deklaration oder in einem Konstruktor der zugehörigen Klasse verändert werden
	\end{itemize}
\end{frame}

\begin{frame}{Schlüsselwörter}
	\lstinputlisting{resources/05_null_exceptions_schluesselwoerter/static_readonly.cs}
\end{frame}

\begin{frame}{Schlüsselwörter}
	\textbf{var}\\
	\begin{itemize}
		\item Kann für Variablen innerhalb von Methoden genutzt werden
		\item Dabei wird der Datentyp der Variable durch den Compiler bestimmt
		\item Syntax:
		\begin{itemize}
			\item \texttt{var \alert{<Variablenname>} = \alert{<Eine Zuweisung>};}
		\end{itemize}
		\item Man sollte es nur nutzen, wenn man durch die rechte Seite erkennt, was für ein Typ zugewiesen wird
	\end{itemize}
	\lstinputlisting{resources/05_null_exceptions_schluesselwoerter/var.cs}
\end{frame}

\end{document}
